% Moved out of dynamic-algebra-potentials.tex on the way to a fully formal
% proof of thm.lax_monoidal. The contents below are correct (modulo their
% own [Proof remains to be supplied or cited.] tags) but introduce
% partial-derivative bookkeeping that the abstract construction handles
% formally; nothing in the main paper depends on them anymore.

%--------------------------------------------------------------------
\section{Supporting corollaries}\label{sec.corollaries}
%--------------------------------------------------------------------

The following two corollaries of the abstract machinery specialized to $\cot\colon\mfd\to\poly$ are used in the proofs below. Both are direct consequences of \cref{prop.cot_monoidal,lem.cot_comonoid,prop.Theta}.

\begin{corollary}[Cotangent chain rule]\label{cor.chain_rule}
For smooth maps $u_i\colon M\to A_i$ ($i=1,\ldots,k$) and $\psi\colon A_1\times\cdots\times A_k\to N$, setting $h\coloneqq\psi\circ\langle u_1,\ldots,u_k\rangle$, we have $\cot{h} = \cot{\psi}\circ(\cot{u_1}\otimes\cdots\otimes\cot{u_k})\circ\cot{\Delta_M^{(k)}}$ in $\poly$. At position $m\in M$ and direction $\sigma\in T^*_{h(m)}N$,
\[
(dh_m)^*\sigma = \sum_{i=1}^k (du_{i,m})^*(\partial_i\psi)^*\sigma: T^*_m M,
\]
where partial derivatives of $\psi$ are evaluated at $(u_1(m),\ldots,u_k(m))$.
\end{corollary}

\begin{proof}
\emph{[Proof remains to be supplied or cited.]} Apply \cref{prop.cot_monoidal} and \cref{lem.cot_comonoid} to factor $\cot h$ through the iterated diagonal, then read the direction-action component-by-component.
\end{proof}

\begin{corollary}[Monoid of covector fields]\label{cor.covector_monoid}
For every manifold $N$, $\ihom{\cot N,\yon}$ is a commutative monoid in $(\poly,\yon,\otimes)$. Its positions correspond to smooth covector fields on $N$ via $\varphi\mapsto\bigl(n\mapsto\bk{\varphi}_n(*)\bigr)$. The $k$-fold multiplication $\mu_N^{(k)}\colon\ihom{\cot N,\yon}^{\otimes k}\to\ihom{\cot N,\yon}$ is pointwise addition:
\[
\mu_N^{(k)}(\alpha_1,\ldots,\alpha_k)(n) = \alpha_1(n)+\cdots+\alpha_k(n).
\]
\end{corollary}

\begin{proof}
\emph{[Proof remains to be supplied or cited.]} The comonoid structure on $\cot{N}$ from \cref{lem.cot_comonoid} dualizes via closedness to a commutative monoid structure on $\ihom{\cot{N},\yon}$. At each position $n\in N$, the comultiplication $\cot{\Delta_N^{(k)}}$ is the $k$-fold fiberwise addition on $T^*_n N$, and its adjoint yields pointwise addition of covector fields.
\end{proof}

%--------------------------------------------------------------------
\section{Hamilton-like equations}\label{sec.hamilton}
%--------------------------------------------------------------------

\subsection{Textbook Hamilton's equations}\label{sec.textbook_hamilton}

We recall Hamilton's equations in the three forms used below; see Marsden--Ratiu~\cite[\S2.4 (canonical form, Prop.~2.4.4), \S7.7 eq.~(7.7.7) (simple mechanical), \S7.8 eq.~(7.8.14) and Thm.~7.8.6 (Lagrange--d'Alembert with external forces)]{marsden1999introduction} for a systematic treatment.

Let $Q$ be a smooth manifold (the \emph{configuration space}) and $H\colon T^*Q\to\rr$ a smooth function (the \emph{Hamiltonian}). In canonical coordinates $(q^i, p_i)$ on $T^*Q$, Hamilton's equations govern the time evolution of a state $(q(t), p(t)): T^*Q$:
\begin{equation}\label{eqn.hamilton_general}
\dot q^i \;=\; \frac{\partial H}{\partial p_i},
\qquad
\dot p_i \;=\; -\frac{\partial H}{\partial q^i}.
\end{equation}
For a \emph{simple mechanical system} with kinetic term $\tfrac{1}{2}p^\top M^{-1}p$ (where $M$ is the inertia tensor, a positive-definite symmetric matrix) and potential $V\colon Q\to\rr$, the Hamiltonian is $H(q,p)=\tfrac{1}{2}p^\top M^{-1}p+V(q)$ and \eqref{eqn.hamilton_general} reduces to
\begin{equation}\label{eqn.hamilton_mechanical}
\dot q \;=\; M^{-1}p,
\qquad
\dot p \;=\; -\nabla V(q).
\end{equation}
When the system is coupled to an environment that can exchange momentum with it, the momentum equation acquires an \emph{external generalized force} $F_\mathrm{ext}(q,p,t)$:
\begin{equation}\label{eqn.hamilton_forced}
\dot q \;=\; M^{-1}p,
\qquad
\dot p \;=\; -\nabla V(q) + F_\mathrm{ext}.
\end{equation}

\subsection{The reduction theorem}\label{sec.reduction}

\begin{theorem}[Reduction]\label{thm.reduction}
Let $f=(P,\outp f,\inpt f,V)\colon\lensob L\to\lensob M$ be a morphism in $\potlens$. By \cref{cor.covector_monoid}, a position of $\Phi_1\lensob M=\cot{\outp M}\otimes\ihom{\cot{\inpt M},\yon}$ is a pair $(m,\beta)$ with $m:\outp M$ and $\beta:\Omega^1(\inpt M)$, and a direction there is a pair $(\nu,m'): T^*_m\outp M\times\inpt M$. At state $(u,\pi): T^*P$ on input $(\ell,\alpha):\Phi_1\lensob L(1)$:
\begin{enumerate}
\item The action of $\Phi(f)$ produces output $(m,\beta):\Phi_1\lensob M(1)$ with
\[
m = \outp f(u,\ell),
\qquad
\beta(m') = (\partial_{\inpt M}\inpt f)^*\alpha\bigl(\inpt f(u,\ell,m')\bigr) - d_{\inpt M}V(u,\ell,m').
\]
\item On received direction $(\nu,m')$, the direction-action on the state slot $T^*_{(u,\pi)}(T^*P)\cong P^*\oplus P$ gives
\[
\dot u = \sharp_P(\pi),
\qquad
\dot\pi = -d_PV + (\partial_P\outp f)^*\nu + (\partial_P\inpt f)^*\alpha\bigl(\inpt f(u,\ell,m')\bigr).
\]
\item The feedback on the input slot is $(l^\mathrm{fb},\mu^\mathrm{fb}):\inpt L\times T^*_\ell\outp L$ with
\[
l^\mathrm{fb} = \inpt f(u,\ell,m'),
\qquad
\mu^\mathrm{fb} = -d_{\outp L}V + (\partial_{\outp L}\outp f)^*\nu + (\partial_{\outp L}\inpt f)^*\alpha\bigl(\inpt f(u,\ell,m')\bigr).
\]
\end{enumerate}
All partials are evaluated at $(u,\ell,m')$.
\end{theorem}

\begin{proof}
\emph{[Proof remains to be supplied or cited.]} The state-slot dynamics specialize \cref{prop.hamilton} to the state-dependent Hamiltonian of \cref{rmk.hamilton_correspondence}: $\dot u=\sharp_P(\pi)$ is $\partial_p H$ for the kinetic-plus-potential Hamiltonian, supplied by the pairing projection $\rho_P$ (\cref{def.pairing_proj}), whose $P$-component is constant in the input direction. The three force-like terms in $\dot\pi$, $\mu^\mathrm{fb}$, and $\beta$ arise from \cref{cor.chain_rule} applied to the three-fold codiagonal on $P\times\outp L\times\inpt M$ along the branches $V$, $\outp f$, $\inpt f$; their sum is the pointwise addition of \cref{cor.covector_monoid}.
\end{proof}

\begin{remark}[Correspondence with textbook Hamilton's equations]\label{rmk.hamilton_correspondence}
With state-dependent Hamiltonian $H(u,\pi)\coloneqq\tfrac{1}{2}\langle\pi,\sharp_P(\pi)\rangle+V(u,\ell,m')$ (input $(\ell,\alpha)$ and received direction $(\nu,m')$ held fixed for the flow), \cref{thm.reduction}'s state update is \cref{prop.hamilton} augmented by the external generalized forces $(\partial_P\outp f)^*\nu+(\partial_P\inpt f)^*\alpha(\inpt f(u,\ell,m'))$, matching the forced form \eqref{eqn.hamilton_forced} with $\sharp_P$ playing the role of the inverse inertia $M^{-1}$. The forces are transmitted through the lens's forward map $\outp f$ (carrying back-force $\nu$ from the target) and its backward map $\inpt f$ (carrying covector field $\alpha$ from the source); a closed system ($\nu=0$, $\alpha=0$) collapses to \eqref{eqn.hamilton_mechanical}.\qedhere
\end{remark}
